\documentclass[fleqn]{article} 
\usepackage{titlesec}
\usepackage{hyperref}
\usepackage{url}
\usepackage{amsmath}
\usepackage{algorithm}
\usepackage[noend]{algpseudocode}
\usepackage{graphicx}
\usepackage{amssymb}
\usepackage{eqnarray}
\usepackage{multicol}
\usepackage{enumitem}
\usepackage{indentfirst}
%\usepackage[utf8]{inputenc}
%\usepackage[T1]{fontenc}
\setlength\columnsep{20pt}
\usepackage[margin={.75in,1.25in}]{geometry}

% Section spacing and font setup
\titlespacing{\section}{0pt}{10pt}{10pt}
\titleformat*{\section}{\large\bfseries}

%Bibliography setup
\usepackage[backend=biber, style=authoryear]{biblatex}
\DeclareNameAlias{sortname}{family-given}
\addbibresource{references.bib}

\title{\huge Planetary Obliquity and Stellar Structure \\ \Large ASTR-154 Final Project}
\date{June 7, 2026}

\author{Babak Aryan\\Justin Kan}

\newcommand{\fix}{\marginpar{FIX}}
\newcommand{\new}{\marginpar{NEW}}

%\nipsfinalcopy % Uncomment for camera-ready version

\begin{document}
	\maketitle

	\section*{ \centering \textbf{ABSTRACT}}
	
	Planetary obliquity, a planet’s axial tilt with respect to its plane of orbit, is the primary driver of seasons and climate cycles for an earth-like planet. An obliquity close to $0^\circ$, for example, eliminates seasonal changes; locations closer to the equator remain hot, while the poles stay perpetually frozen. A large obliquity ($> \; \sim\!45^\circ$), on the other hand, creates large variations in climate across different altitudes on the planet resulting in the hemispheres cycling between very hot summers and brutally cold winters over an orbital year.  According to (\cite{Vervoort}), one important condition to increase the probability of “the habitability of Earth-like planets” is for the mean obliquity to be “sufficiently low to maintain temperate temperatures over large parts of its surface throughout the orbital year.”  Planetary obliquity, therefore, should be considered as one of the essential parameters when searching for habitable exoplanets. Recent space telescope missions such as Kepler/K2, Tess and CoRoT have contributed to the discovery of $\sim40,000$ exoplanets cataloged in NASA Exoplanet Archive. Obliquity data, however, is available for only a handful ($274$) of these planets. Therefore, an efficient search for "habitable" exoplanets necessitates a better understanding of planetary obliquity population in order to better identify suitable candidates. Recent research on the evolution of planetary obliquity has shown the effect of non-single ("binary") stellar host systems on mean equilibrium value ($\epsilon_m$) as well as long term variation ($\Delta\epsilon$) of planetary obliquity (\cite{Huang}), (\cite{Quarles}).	In this paper, we explore the planetary obliquity populations as a function of the planet's host system's "binarity"\footnote{Throughout this paper "binarity" refers to the condition of a system containing more than 1 host star.}. We show that obliquity population for planets in single stellar host systems differ from the population for those in "binary" systems. In order to divide our exoplanet sample set into those in single star systems vs. those in binary systems, we rely on the number of listed stars in the data set as well as the test statistics RUWE (Re-normalized Unit Weight Error) which is a strong indicator of the binarity of a system (\cite{RUWE_Ginard}), (\cite{RUWE_Logan}). We show that, given a grouping of exoplanets based on both the number of stars listed and the RUWE number, yields obliquity populations which are distinctly different with a $p$ value of $< 0.03$. Given the high correlation between RUWE and binarity, and how binarity informs the planetary obliquity population, we also explore the correlation between RUWE and planetary obliquity. We show that there is a small positive correlation between RUWE and planetary obliquity for the population in single stellar systems while the correlation is one of a small negative one for the population in binary systems.
	

	\begin{multicols}{2}
		\section*{Motivation}
		Exoplanet detection is a growing field of study in the Astronomy and Astrophysics field. To date there are 6,298 confirmed exoplanets (\cite{ExoArchive}).
		\begin{figure*}[t]
			\centering
			\includegraphics[width=0.9\linewidth]{../results/confirmed_exoplanets}
			\caption{Confirmed Exoplanet Discoveries by year and method. Data from NASA Exoplanet Archive (NASA Exoplanet Archive, 2026). This is a histogram plot with the number of exoplanets discovered in the y-axis plotted by discovery year in the x-axis. The method of discovery is also described by different colors in the plot with the Transit method in purple.}
			\label{fig:confirmedexoplanets}
		\end{figure*}
		A majority of these are discovered via the transit method\footnote{Types of detection methods: https://sci.esa.int/web/exoplanets/-/60655-detection-methods}. Obliquity, or spin-orbit axis, is an angular measure of a planet’s axial tilt with respect to its plane of orbit. Obliquities vary widely within any given system. The obliquity of any given planet has implications regarding habitability. For climate, an obliquity close to 0o eliminates seasonal changes, locations closer to the equator remain hot, and the poles remain perpetually frozen. A large obliquity ($> \; \sim\!45^\circ$) creates large variations in climate across different altitudes on the planet. Throughout an orbital year, the hemispheres cycle between very hot summers and brutally cold winters. Habitability is tied to climate and therefore to obliquity, so one of the conditions to increase the probability of “the habitability of Earth-like planets” is for the mean obliquity to be “sufficiently low to maintain temperate temperatures over large parts of its surface throughout the orbital year.” (\cite{Vervoort}). In addition to the mean equilibrium obliquity, variations in obliquity also affect climate and have an impact where an exoplanet’s potential to host life may be correlated to large variations (\cite{Spiegel}).
		Research in planetary obliquity evolution also shows that binarity can affect the mean equilibrium value as well as long term variation of planetary obliquity (\cite{Huang}), (\cite{Quarles}). For example, Earth’s obliquity varies between $22.1^\circ$ and $24.5^\circ$ over ~41,000 yr resulting in a very gradual and mild climate variation of time. By comparison, simulations have shown that exoplanets in binary systems can range in mean obliquities between 30º and $60^\circ$ with larger time variations reaching maximum obliquities over $120^\circ$.
		
		\section*{Data}
		Two publicly available data sets were used for this project. NASA Exoplanet Archive and Gaia DR3. Vetted data in the Exoplanet Archive is continually added by a team of astronomers and are linked back to the original literature (\cite{ExoArchive}). Gaia’s full dataset from the Gaia Mission is released periodically. The most current release, DR3, was made available in 2022 and contained over 1.8 billion objects with 1.4 billion sources containing full astrometry (\cite{GaiaDR3}).
		In NASA Exoplanet Archive three columns were queried for pre-processing: gaia\_dr3\_source\_id, sy\_snum, and pl\_projobliq. To analyze data on objects with recorded data these columns were filtered with "not null" in the query yielding 263 records. This was then downloaded as a csv file and read in as an Astropy Table. For Gaia the column gaia\_dr3\_source\_id was cross-referenced with the same value from the Exoplanet Archive, yielding 216 records. This subset was then downloaded as a csv file and read to an Astropy table.
		The RUWE value from Gaia DR3 is calculated by utilizing the astrometric instrument in the Gaia mission. It is comprised of two identical three-mirror anastigmatic (TMA) telescopes with apertures 1.45 m x 0.50 m. These telescopes are used in conjunction with 14 CCD\footnote{Charge-coupled devices, model: e2v technologies CCD91-72 (\cite{Bruijne})}s and an additional area of 62 CCDs to develop an astrometric field in the 330-1050 nm range. Full details on instrumentation detailed in sections 3.3.1and 3.3.5 (\cite{GaiaDr3-2016}).
\		
		\section*{Methods}
		Describe your analysis procedure. Why did you choose the techniques you used? If you used a technique beyond what we learned this quarter, teach it to the class. If there were other techniques you considered but chose not to use, explain why. 
		
		\section*{Results}
		What did you learn by applying your methods to your data? How does this help you answer the question(s) posed in the motivation? Include a discussion of the uncertainty or statistical significance of your results.
		
		\section*{Next Steps}
		If you were to continue along this line of research, how could you expand on your project to answer the question(s) better or to answer related questions? What data and techniques would you need? If your project didn't work out as expected, what do you think went wrong, and how could you improve?
	\end{multicols}
	
	\pagebreak
	
	\printbibliography

\end{document}

